\documentclass[11pt]{article}
\usepackage[utf8]{inputenc}
\usepackage[T1]{fontenc}
\usepackage[margin=1in]{geometry}
\usepackage{hyperref}
\usepackage{parskip}
\usepackage{enumitem}

% Cover letter for JMLR submission of the VSP negative-result paper.
% Modeled on jmlr_cover_letter.tex (softmax paper). Keep the two consistent.
% Before submission: confirm the suggested Action Editors and Reviewers against
% the current JMLR Action Editor list (names below are lane-appropriate but the
% live roster changes).

\begin{document}

\thispagestyle{empty}

\noindent Nikita Gorshkov\\
Senefelderstr.\ 29a, 10437 Berlin, Germany\\
\texttt{ngorshkov@proton.me}\\[1em]
\noindent\today\\[1em]

\noindent To the Editors of the \emph{Journal of Machine Learning Research},\\

I am submitting the manuscript \emph{Grounded Token Bundles Separate Word Senses
but Do Not Help a Language Model: A Negative Result} for consideration for
publication in JMLR.

\paragraph{Summary.}
The paper tests whether a grounded token representation, one bundling a word
sense's visual (V), semantic (S), and physical (P) features into a single vector,
improves word-sense disambiguation in a small autoregressive language model. It
separates the hypothesis into two claims. The representation claim holds: on
curated polysemous words the V+S+P bundle attains a sense-separation of $0.37$
(max-aggregated $1-\cos$) where the text block alone attains $0.00$ by
construction, and a scrambled-bundle control confirms the signal is genuine (a
true bundle initializes a model $+10.5$ points better than a permuted one at low
compute, $p = 0.007$). The transfer claim fails: across two delivery mechanisms
of the same bundle (embedding initialization and inference-time reranking) and
three benchmarks, the grounding produces no improvement a trained model retains.
The full-compute initialization deficit reproduces across three seeds ($-3.5$
points, $95\%$ CI $[-4.9, -2.1]$), and by full compute a plain random start
outperforms the grounded one. The paper's most transferable contribution is
methodological: seed reproduction and held-out hyperparameter selection converted
three separately promising effects ($+5.7$, $+2.9$, $+8.3$ points) into a single
bounded negative. We report the result as a bound, not a proof of absence, and are
explicit about the underpowered individual tests.

\paragraph{Fit for JMLR.}
JMLR explicitly welcomes negative and reproducibility-focused results of clear
methodological value. This submission is such a case: a clean negative with an
explicit diagnosis, honest power analysis, adversarial controls, and released
code and result files for every reported number.

\section*{1.~Overlapping or prior publications}
\begin{itemize}[leftmargin=1.4em]
  \item This manuscript has \textbf{not} been published in, or simultaneously
    submitted to, any other journal or conference.
  \item No preprint of this manuscript has been posted (no arXiv or other
    preprint server). This is the first public release of the work.
  \item The author has separately submitted an unrelated theoretical manuscript
    to JMLR (on the expressiveness of alpha-compositing versus softmax attention).
    The two share no results, figures, or text; both arise from the same code
    repository, which is cited.
\end{itemize}

\section*{2.~Co-author awareness and consent}
Single-author submission (Nikita Gorshkov). No co-authors to notify.

\section*{3.~Conflicts of interest}
The author declares \textbf{no conflicts of interest} with the suggested action
editors or reviewers (no family/close-friend, advisor/advisee, or collaboration
within the past three years), and \textbf{no funding and no competing interests}
for the prior 36 months.

\section*{4.~Suggested Action Editors}
Drawn from the current JMLR Action Editor list. The paper's lanes are empirical
NLP (word-sense disambiguation), multimodal/grounded representation learning, and
reproducibility of negative results; the suggestions below are chosen for those
lanes. \emph{Please confirm against the live AE roster, which changes over time.}
\begin{enumerate}[leftmargin=1.6em]
  \item \textbf{Ivan Titov} --- natural language processing and representation
    learning; fit for the WSD and sense-representation core.
  \item \textbf{Trevor Cohn} --- statistical NLP and language modeling; fit for
    the language-model transfer evaluation.
  \item \textbf{Kevin Gimpel} --- NLP, representation learning, embeddings; fit
    for the sense-embedding and disambiguation framing.
\end{enumerate}

\section*{5.~Suggested Reviewers}
The author has no collaboration with any of these. Affiliations are given for
identification; contact addresses can be confirmed on request.
\begin{enumerate}[leftmargin=1.6em]
  \item \textbf{Roberto Navigli} --- Sapienza University of Rome; author of the
    canonical word-sense-disambiguation survey and long-running WSD research.
  \item \textbf{Douwe Kiela} --- Stanford University / Contextual AI; multimodal
    and grounded representation learning.
  \item \textbf{Elia Bruni} --- multimodal distributional semantics
    (``Multimodal Distributional Semantics''); fit for the grounding claim.
  \item \textbf{Jesse Dodge} --- Allen Institute for AI; reproducibility and
    honest reporting of experimental results (cited in the paper). (optional)
\end{enumerate}

\section*{6.~Keywords}
word-sense disambiguation; grounded representations; multimodal embeddings;
language models; negative results; reproducibility.

\bigskip
\noindent Thank you for considering this submission.

\bigskip
\noindent Sincerely,\\[0.5em]
Nikita Gorshkov (corresponding author)\\
Senefelderstr.\ 29a, 10437 Berlin, Germany\\
\texttt{ngorshkov@proton.me}

\end{document}
